\chapter{Comunicaciones}
\label{chap:comunicaciones}

El sistema usa cuatro enlaces (tabla~\ref{tab:enlaces}). Este capítulo describe el protocolo binario
propio entre el microcontrolador y el ordenador, el uso de Modbus RTU con el accionamiento y la
interfaz entre el servidor y el navegador. La especificación byte a byte está en el
anexo~\ref{anx:protocolo}.

\begin{table}[H]
\centering
\caption{Enlaces de comunicación del sistema.}
\label{tab:enlaces}
\small
\begin{tabularx}{\textwidth}{@{}lL{0.27\textwidth}XX@{}}
\toprule
Enlace & Extremos & Capa física & Protocolo \\
\midrule
USB & PC $\leftrightarrow$ Pico & USB CDC (puerto serie virtual) & Binario propio + CRC-8 \\
RS-485 & Pico $\leftrightarrow$ accionamiento & UART PIO, 115200 8N1 & Modbus RTU \\
I\textsuperscript{2}C1 & Pico $\leftrightarrow$ ZSC31014 & \SI{400}{\kilo\hertz} (\SI{100}{\kilo\hertz} en modo comando) & Propio del chip \\
Red local & Navegador $\leftrightarrow$ servidor & TCP, puerto 8080 & HTTP + WebSocket (JSON) \\
\bottomrule
\end{tabularx}
\end{table}

\section{Protocolo binario entre el microcontrolador y el ordenador}

\subsection{Formato de trama}

Todas las tramas, en ambos sentidos, tienen el mismo formato (figura~\ref{fig:trama}): dos bytes de
sincronía \codigo{0xAA 0x55}, un byte de longitud, la carga útil y un CRC-8 calculado sobre la carga
útil. El primer byte de la carga útil identifica el tipo de paquete o la orden. Los enteros
multibyte van en orden \emph{little-endian}.

\begin{figure}[H]
\centering
\begin{bytefield}[bitwidth=2.4em]{12}
  \bitbox{2}{\small 0xAA} & \bitbox{2}{\small 0x55} & \bitbox{2}{\small LEN} &
  \bitbox{1}{\small ID} & \bitbox{3}{\small datos (LEN$-$1)} & \bitbox{2}{\small CRC-8}
\end{bytefield}
\caption{Formato común de trama. El CRC-8 cubre los LEN bytes de la carga útil (ID + datos).}
\label{fig:trama}
\end{figure}

El CRC usa el polinomio $x^8+x^2+x+1$ (\codigo{0x07}), valor inicial \codigo{0x00}, sin reflexión ni
XOR final (la variante conocida como CRC-8/SMBUS). Se implementa bit a bit en el firmware y en
Python con el mismo algoritmo. Para tramas cortas como éstas detecta todos los errores de uno y dos
bits y todas las ráfagas de hasta 8 bits.

\subsection{Recepción y resincronización}

El firmware decodifica las órdenes con una máquina de cuatro estados (\codigo{RX\_SYNC0},
\codigo{RX\_SYNC1}, \codigo{RX\_LEN}, \codigo{RX\_PAYLOAD}). Una longitud nula o mayor de 60 devuelve
el decodificador a la búsqueda de sincronía, y una trama con CRC incorrecto se descarta sin
ejecutarse. Si se pierden bytes, el receptor vuelve a encontrar el patrón \codigo{AA 55} en la
siguiente trama. El servidor hace lo mismo con la telemetría, con un tiempo de espera de
\SI{100}{\milli\second} por lectura.

Para depurar desde un terminal, el firmware acepta también órdenes de una letra (\codigo{i},
\codigo{a}, \codigo{b}, \codigo{x}, \codigo{h}, \codigo{q}), pero \emph{sólo} cuando el decodificador
está esperando el primer byte de sincronía. Antes se comprobaban en cualquier estado, y eso
provocaba un error difícil de ver: la orden de programar la célula con ganancia 7 y cero 8 genera la
trama \codigo{AA 55 03 09 07 08 69}, cuyo CRC \codigo{0x69} coincide con el carácter \codigo{i}. El
atajo se comía el último byte, la orden se perdía y además se reinicializaba el servo. Con ganancia
6 el CRC vale \codigo{0x7C} y la misma orden funcionaba. Ambos valores se han comprobado con la
implementación del CRC.

\subsection{Paquetes del microcontrolador al ordenador}

\begin{table}[H]
\centering
\caption{Paquetes de telemetría y estado.}
\label{tab:paquetes}
\begin{tabular}{@{}clll@{}}
\toprule
ID & Nombre & Tamaño de carga útil & Cuándo se envía \\
\midrule
\codigo{0x02} & Telemetría por lotes & $14 + 14n$ B, $n\le 8$ & cada \SI{200}{\milli\second} o con 8 muestras \\
\codigo{0x03} & Estado de \emph{homing} & 11 B & junto al lote, mientras hay \emph{homing} \\
\codigo{0x04} & Resultado de la célula & 13 B & una vez por consulta o programación \\
\codigo{0x05} & Estado del ciclo & 18 B & junto al lote, mientras hay ciclo \\
\bottomrule
\end{tabular}
\end{table}

\paragraph{Telemetría por lotes.} Cada vuelta del lazo añade al búfer una muestra de 14 bytes:
desfase en milisegundos respecto al instante base del lote, velocidad, par, consigna, fuerza tarada,
estado de los finales de carrera, estado de la máquina y tiempo de trabajo de la vuelta. Los datos
comunes (instante base de 32 bits, lectura bruta de la célula y su estado, indicadores, código de
error, dirección del servo y códigos de resultado de las dos lecturas Modbus) viajan una sola vez en
una cabecera de 14 bytes.

Agrupar las muestras reduce las escrituras por USB a cinco por segundo sin perder resolución
temporal, porque cada muestra conserva su instante exacto. Con el lazo de \SI{50}{\milli\second} cada
lote lleva típicamente cuatro muestras: 70 bytes de carga útil, 74 de trama y unos
\SI{370}{\byte\per\second}, muy por debajo de la capacidad del enlace. El precio es un retardo de
visualización de hasta \SI{200}{\milli\second} y que la lectura bruta de la célula sólo se refresque
una vez por lote (sección~\ref{sec:res-limitaciones}).

\paragraph{Estado del ciclo.} El paquete \codigo{0x05} creció de 14 a 18 bytes al reorganizar el
ciclo: además de la fase, las repeticiones y la posición, transporta \codigo{cycPosB}, la posición
en la que ha quedado fijado el punto B. Sin ese dato no se puede saber desde fuera contra qué está
midiendo el ciclo un recorrido parcial. Las fases son ahora seis: reposo, colocándose en B, hacia A,
volviendo a B, completado y abortado.

Para que un firmware y un servidor desparejados no se queden mudos, el servidor acepta también la
trama de 14 bytes de la versión anterior y asume \codigo{cycPosB = 0}. Esa compatibilidad se añadió
precisamente porque un desajuste de versiones despistó durante la depuración del ciclo: el firmware
enviaba una trama que el servidor rechazaba entera, y la interfaz no mostraba ninguna fase.

\paragraph{Modelo pregunta-respuesta de la célula.} Cada orden de consulta, programación o
mantenimiento de la célula produce exactamente un paquete \codigo{0x04}, también cuando se rechaza.
El paquete incluye el código de la orden a la que responde, para que la interfaz nunca se quede
esperando, y lleva información de diagnóstico: el primer byte que devolvió el chip, si el pin de
reinicio corta de verdad la alimentación y en qué intento del barrido de retardos se consiguió
entrar en modo comando (sección~\ref{sec:prob-zsc}).

\subsection{Órdenes del ordenador al microcontrolador}

\begin{longtable}{@{}clL{0.30\textwidth}L{0.33\textwidth}@{}}
\caption{Órdenes aceptadas por el firmware.}\label{tab:ordenes}\\
\toprule
Código & Nombre & Datos & Efecto \\
\midrule
\endfirsthead
\toprule
Código & Nombre & Datos & Efecto \\
\midrule
\endhead
\codigo{0x01} & INIT & [dirección Modbus] opcional & Reinicia la posición, invalida el recorrido medido y configura y habilita el servo. \\
\codigo{0x02} & STOP & -- & Consigna a cero; aborta \emph{homing} y ciclo. \\
\codigo{0x03} & SET\_PARAM & id (1 B), valor int16 & Modifica un parámetro (tabla~\ref{tab:parametros}). \\
\codigo{0x04} & MOVE\_A & -- & Movimiento manual hacia A (arriba). \\
\codigo{0x05} & MOVE\_B & -- & Movimiento manual hacia B (abajo). \\
\codigo{0x06} & SHUTDOWN & -- & Deshabilita el accionamiento. \\
\codigo{0x07} & SCAN & -- & Búsqueda bloqueante del esclavo (1--247). \\
\codigo{0x08} & HOME & -- & Inicia el \emph{homing}. \\
\codigo{0x09} & LC\_PROGRAM & B\_Config, ZMDI\_Config1 (uint16) & Graba la configuración de la célula (sólo parado). \\
\codigo{0x0A} & LC\_QUERY & -- & Lee la configuración de la célula (sólo parado). \\
\codigo{0x0B} & LC\_HOLD & modo de aparcado (1 B) & Mantiene la célula sin tensión \SI{8}{\second} para medirla. \\
\codigo{0x0C} & CYCLE\_START & modo, recorrido (centi-rev), velocidad, repeticiones & Inicia el ciclo B$\rightarrow$A$\rightarrow$B, empezando por la colocación en B. \\
\codigo{0x0D} & CYCLE\_STOP & -- & Detiene el ciclo donde esté. \\
\bottomrule
\end{longtable}

\begin{table}[H]
\centering
\caption{Parámetros configurables con SET\_PARAM.}
\label{tab:parametros}
\begin{tabular}{@{}cllr@{}}
\toprule
Id & Parámetro & Unidad & Defecto \\
\midrule
1 & Velocidad de movimiento manual & rpm & 10 \\
2 & Límite de fuerza (0 = desactivado) & cuentas & 0 \\
3 & Tara de la célula & cuentas & 8210 \\
4 & Aceleración (C03.22) & u. del accionamiento & 5 \\
5 & Deceleración (C03.24) & u. del accionamiento & 5 \\
6 & Rigidez (C00.05) & nivel & 10 \\
7 & Calibración de fuerza (sólo informativa) & cuentas/N & 0 \\
8 & Desplazamiento del origen tras el \emph{homing} & centi-rev & 0 \\
9 & Velocidad de \emph{homing} & rpm & 30 \\
\bottomrule
\end{tabular}
\end{table}

\section{Modbus RTU con el accionamiento}

\subsection{Registros utilizados}

El firmware sólo usa las funciones 0x03 y 0x06 sobre un conjunto reducido de registros
(anexo~\ref{anx:modbus}). Los parámetros del accionamiento se nombran \codigo{Cgg.ii}, y todas las
direcciones cumplen la regla
\begin{equation}
  \text{dirección} = 256\cdot gg + \text{0x}ii,
\end{equation}
es decir, el índice se interpreta en hexadecimal. Por ejemplo, C04.11 (habilitación) es
$1024 + 17 = 1041$ y C03.22 (aceleración) es $768 + 34 = 802$. Las magnitudes de supervisión están en
el grupo \codigo{0x40}: velocidad en 16385 (\codigo{0x4001}) y par en 16387 (\codigo{0x4003}). Tener
esta regla evita errores al añadir registros a partir del manual.

\subsection{Secuencia de configuración}

Al recibir INIT, el firmware configura el accionamiento así: deshabilitar (C04.11 = 0), seleccionar
el modo velocidad (C00.00 = 1), poner la consigna a cero (C03.21 = 0), escribir la rigidez y las
rampas, configurar la resistencia de frenado externa (C00.10 a C00.13) y habilitar (C04.11 = 1). Si
una escritura esencial falla se reintenta la secuencia completa, y tras tres fallos el sistema pasa
a estado de error. Entre escrituras se deja un silencio de \SI{700}{\micro\second}.

\subsection{Transacciones por periodo}

En régimen normal, cada vuelta del lazo realiza tres transacciones: lectura de la velocidad, lectura
del par y escritura de la consigna. Escribir la consigna en cada vuelta, aunque no cambie, sirve
además de supervisión del enlace: si la escritura falla, se marca el servo como desconectado.

A 115200 baudios con formato 8N1, cada byte ocupa \SI{86.8}{\micro\second}. Las tres transacciones
suman 46 bytes, unos \SI{4.0}{\milli\second} de línea. A eso se añaden \SI{0.6}{\milli\second} de
retardos de conmutación del transceptor y \SI{0.7}{\milli\second} de silencio entre las dos lecturas.
El tiempo de trabajo medido del lazo es \SI{10.97}{\milli\second}
(sección~\ref{sec:res-temporizacion}); la diferencia, unos \SI{5.7}{\milli\second}, corresponde
principalmente al tiempo de respuesta del accionamiento, aproximadamente \SI{1.9}{\milli\second} por
transacción.

La guía de implementación de Modbus serie recomienda, por encima de 19200 baudios, silencios fijos
de \SI{1.75}{\milli\second} entre tramas \cite{modbusserial}. El firmware usa
\SI{700}{\micro\second} y el accionamiento lo acepta sin errores, pero con otros esclavos convendría
respetar el valor recomendado.

\section{Interfaz entre el servidor y el navegador}

\subsection{WebSocket}

El servidor mantiene una conexión WebSocket en \codigo{/ws} con cada navegador. En sentido servidor
$\rightarrow$ navegador envía objetos JSON de cinco tipos:
\begin{itemize}
  \item \textbf{muestra de telemetría} (sin campo \codigo{type}): una por muestra del lote, con los
        valores brutos, los nombres de estado y error ya traducidos, la posición integrada, los
        finales de carrera, la frecuencia de muestreo medida y el estado del \emph{homing};
  \item \codigo{home}: fase del \emph{homing}, posición y recorrido medido;
  \item \codigo{cycle}: fase, repeticiones hechas y pedidas, posición, recorrido pedido, posición de
        B y fichero en grabación;
  \item \codigo{lc\_config}: resultado de la consulta o programación de la célula;
  \item \codigo{cmd\_ack}: confirmación de \emph{toda} orden recibida, con indicación de éxito y un
        texto explicativo. Sin ella, una orden que no llegaba al puerto era indistinguible de una
        que sí se había ejecutado.
\end{itemize}

En sentido navegador $\rightarrow$ servidor, las órdenes son objetos
\codigo{\{"cmd": nombre,} \codigo{...parámetros\}}. El servidor las valida, limita los parámetros a
rangos seguros (por ejemplo, velocidad del ciclo entre 0 y \SI{300}{rpm}), construye la trama
binaria y la escribe en el puerto.

\subsection{Servicios HTTP}

\begin{table}[H]
\centering
\caption{Servicios HTTP del servidor.}
\label{tab:http}
\begin{tabular}{@{}lll@{}}
\toprule
Método & Ruta & Función \\
\midrule
GET & \codigo{/} & Monitor en tiempo real \\
GET & \codigo{/ensayos} & Visor de ensayos grabados \\
GET & \codigo{/sim} & Configuración del mecanismo y simulador \\
GET & \codigo{/api/mecanismo} & Geometría actual (JSON) \\
PUT & \codigo{/api/mecanismo} & Guarda la geometría, validada contra rangos \\
GET & \codigo{/cycles} & Lista de los 50 ensayos más recientes (JSON) \\
GET & \codigo{/cycles/\{nombre\}} & Descarga de un ensayo (CSV) \\
DELETE & \codigo{/cycles/\{nombre\}} & Borrado de un ensayo (rechazado si se está grabando) \\
GET & \codigo{/mecanismo.js}, \codigo{/tema.css}, \codigo{/tema.js}, \codigo{/nav.js} & Recursos comunes a las tres páginas \\
\multicolumn{2}{@{}l}{\codigo{/fui} (montaje estático)} & Kit de interfaz (hoja de estilo y tipografías) \\
\bottomrule
\end{tabular}
\end{table}

Las rutas con nombre de fichero reducen el parámetro a su último componente antes de usarlo, de modo
que una petición con \codigo{../} no puede leer ni borrar ficheros fuera de la carpeta de ensayos.
La escritura de la geometría se valida en el servidor contra los mismos rangos que comprueba el
navegador, de forma que una petición construida a mano tampoco puede dejar una configuración
imposible.
